\section{Conclusion}

This paper presented an execution evidence architecture for agentic software systems and instantiated it as a reproducible reference path in an artifact package omitted here for double-blind review. The central claim is that observability alone is not enough: agent runs need structured evidence bundles that make intent, authorization, execution verification, and exported receipts independently reviewable.

The checked-in evaluation harness supports a bounded version of that claim with a deterministic 16-task suite, a stable five-file export contract, a third-party review script, minimal live CrewAI and LangChain baselines, paired same-framework comparisons in both frameworks with and without the evidence layer, four ablation modes, and a small negative-control falsification suite. An appendix-only blinded reviewer workflow is also included as reproducibility scaffolding. Across the current summaries, the full evidence chain preserves review questions that the trace-oriented baseline and the live framework observability baselines do not explicitly retain, the same live runtime becomes reviewable when wrapped by the evidence layer in two framework ecosystems, the architecture shows interpretable degradation under targeted ablations, and the same review contract rejects several malformed or over-claimed bundles derived from previously valid runs.

Future work should validate the architecture against live framework integrations, extend the integrity layer beyond the current bounded checks, and run a real blinded reviewer study. Even in its current form, the reference implementation shows a plausible and reproducible path for engineering execution evidence as a first-class systems layer rather than treating it as a side effect of logging.
